% !TeX root = part02.tex

\documentclass[../final.tex]{subfiles}

\begin{document}

\chapter{Спектр атома щелочного металла}

\rsubtitle{Атомы щелочных металлов}

\section{Одноэлектронная модель щелочного атома}

Рассмотрим многоэлектронный атом с зарядом ядра \(+Ze\), содержащий
\(N\) электронов. Потенциальную энергию такой системы можно записать в виде
\begin{equation*}
    U
    =
    -\sum_{i=1}^{N}\frac{Ze^2}{r_i}
    +
    \frac{1}{2}
    \sum_{\substack{i,k=1\\ i\neq k}}^{N}
    \frac{e^2}{r_{ik}}.
\end{equation*}

Первое слагаемое описывает кулоновское притяжение электронов к ядру,
а второе --- кулоновское отталкивание электронов друг от друга.

Для атомов щелочных металлов все внутренние электронные оболочки замкнуты,
а на внешней оболочке находится один валентный электрон. Например, для натрия
\begin{equation*}
    \mathrm{Na}:
    1s^2\,2s^2\,2p^6\,3s.
\end{equation*}

Поэтому щелочной атом удобно рассматривать как систему из атомного остова
и одного внешнего электрона.

Обозначим через \(\mathbf r\) радиус-вектор внешнего электрона,
а через \(\mathbf r_k\) --- координаты \(N-1\) внутренних электронов.
Отбрасывая постоянную энергию самого атомного остова, для внешнего электрона
получим гамильтониан
\begin{equation*}
    \hat H
    =
    \frac{\hat{\mathbf p}^{\,2}}{2m}
    -
    \frac{Ze^2}{r}
    +
    \sum_{k=1}^{N-1}
    \frac{e^2}{|\mathbf r-\mathbf r_k|}.
\end{equation*}

% ============================================================
% РИСУНОК:
% Схематическое изображение атома щелочного металла:
% ядро +Ze, внутренние N-1 электронов, образующие атомный остов,
% и внешний электрон на расстоянии r.
% ============================================================

\section{Поле атомного остова}

Рассмотрим внешний электрон на расстоянии, значительно превышающем
характерный размер атомного остова:
\begin{equation*}
    r\gg r_k.
\end{equation*}

Тогда взаимодействие внешнего электрона с каждым внутренним электроном
можно разложить по малому параметру \(r_k/r\):
\begin{align}
    \frac{1}{|\mathbf r-\mathbf r_k|}
    &=
    \frac{1}
    {\sqrt{r^2+r_k^2-2\mathbf r\cdot\mathbf r_k}}
    \\
    &=
    \frac{1}{r}
    \frac{1}
    {\sqrt{
        1+\dfrac{r_k^2}{r^2}
        -2\dfrac{\mathbf r\cdot\mathbf r_k}{r^2}
    }}.
\end{align}

В первом приближении
\begin{equation*}
    \frac{1}{|\mathbf r-\mathbf r_k|}
    \simeq
    \frac{1}{r}
    \left(
        1+
        \frac{\mathbf r\cdot\mathbf r_k}{r^2}
        +\ldots
    \right).
\end{equation*}

Подставляя это разложение в гамильтониан, получаем
\begin{align}
    \hat H
    \simeq
    \frac{\hat p^2}{2m}
    &-
    \frac{Ze^2}{r}
    +
    (N-1)\frac{e^2}{r}
    \nonumber\\
    &+
    \frac{e^2}{r^3}
    \mathbf r\cdot
    \sum_{k=1}^{N-1}\mathbf r_k
    +\ldots
\end{align}

Нулевой член разложения даёт
\begin{align}
    -\frac{Ze^2}{r}
    +(N-1)\frac{e^2}{r}
    &=
    -\frac{(Z-N+1)e^2}{r}.
\end{align}

Введём эффективный заряд атомного остова
\begin{equation*}
    \boxed{
        Z_a=Z-N+1
    }.
\end{equation*}

Для нейтрального щелочного атома
\begin{equation*}
    N=Z,
\end{equation*}
поэтому
\begin{equation*}
    \boxed{
        Z_a=1
    }.
\end{equation*}

Таким образом, на больших расстояниях внешний электрон видит атомный остов
как объект с единичным положительным зарядом. В нулевом приближении
\begin{equation*}
    U_0(r)
    =
    -\frac{Z_a e^2}{r}.
\end{equation*}

\subsection{Поляризация атомного остова}

Следующая поправка связана с поляризацией атомного остова полем внешнего
электрона. В используемом приближении её запишем как
\begin{equation*}
    U_1(r)
    =
    -C_1\frac{Z_a e^2}{r^2},
\end{equation*}
где \(C_1\) характеризует величину поляризационного эффекта.

В результате эффективная потенциальная энергия внешнего электрона имеет вид
\begin{equation*}
    \boxed{
        U(r)
        =
        -\frac{Z_a e^2}{r}
        -
        C_1\frac{Z_a e^2}{r^2}
    }.
\end{equation*}

Именно наличие второго слагаемого отличает движение внешнего электрона
щелочного атома от движения электрона в чистом кулоновском поле.

\section{Уравнение Шрёдингера для внешнего электрона}

Запишем стационарное уравнение Шрёдингера
\begin{equation*}
    \hat H\psi=E\psi.
\end{equation*}

Гамильтониан имеет вид
\begin{equation*}
    \hat H
    =
    -\frac{\hbar^2}{2m}\Delta
    -
    \frac{Z_a e^2}{r}
    -
    C_1\frac{Z_a e^2}{r^2}.
\end{equation*}

В сферических координатах
\begin{equation*}
    \Delta
    =
    \frac{1}{r^2}
    \frac{\partial}{\partial r}
    \left(
        r^2\frac{\partial}{\partial r}
    \right)
    +
    \frac{1}{r^2}\Delta_\Omega,
\end{equation*}
где \(\Delta_\Omega\) --- угловая часть оператора Лапласа.

Поэтому
\begin{align}
    -\frac{\hbar^2}{2mr^2}
    \frac{\partial}{\partial r}
    \left(
        r^2\frac{\partial\psi}{\partial r}
    \right)
    &-
    \frac{\hbar^2}{2mr^2}
    \Delta_\Omega\psi
    -
    \frac{Z_a e^2}{r}\psi
    \nonumber\\
    &-
    C_1\frac{Z_a e^2}{r^2}\psi
    =
    E\psi.
\end{align}

Так как потенциал центрально-симметричен, разделим переменные:
\begin{equation*}
    \psi(r,\theta,\varphi)
    =
    R(r)Y_{\ell m}(\theta,\varphi).
\end{equation*}

Для сферических функций
\begin{equation*}
    \Delta_\Omega Y_{\ell m}
    =
    -\ell(\ell+1)Y_{\ell m}.
\end{equation*}

Следовательно, радиальное уравнение принимает вид
\begin{align}
    -\frac{\hbar^2}{2mr^2}
    \frac{d}{dr}
    \left(
        r^2\frac{dR}{dr}
    \right)
    &+
    \frac{\hbar^2\ell(\ell+1)}{2mr^2}R
    -
    C_1\frac{Z_a e^2}{r^2}R
    \nonumber\\
    &-
    \frac{Z_a e^2}{r}R
    =
    ER.
\end{align}

Соберём вместе слагаемые, пропорциональные \(1/r^2\):
\begin{align}
    -\frac{\hbar^2}{2mr^2}
    \frac{d}{dr}
    \left(
        r^2\frac{dR}{dr}
    \right)
    &+
    \frac{\hbar^2}{2mr^2}
    \left[
        \ell(\ell+1)
        -
        C_1\frac{2mZ_a e^2}{\hbar^2}
    \right]R
    \nonumber\\
    &-
    \frac{Z_a e^2}{r}R
    =
    ER.
\end{align}

\section{Квантовый дефект}

Чтобы привести радиальное уравнение к водородоподобному виду,
введём эффективное орбитальное квантовое число \(\tilde\ell\):
\begin{equation*}
    \boxed{
        \tilde\ell(\tilde\ell+1)
        =
        \ell(\ell+1)
        -
        C_1\frac{2mZ_a e^2}{\hbar^2}
    }.
\end{equation*}

Это квадратное уравнение относительно \(\tilde\ell\):
\begin{equation*}
    \tilde\ell^2+\tilde\ell
    -
    \ell(\ell+1)
    +
    C_1\frac{2mZ_a e^2}{\hbar^2}
    =
    0.
\end{equation*}

Отсюда
\begin{equation*}
    \tilde\ell
    =
    -\frac12
    +
    \frac12
    \sqrt{
        1+4\ell(\ell+1)
        -
        C_1\frac{8mZ_a e^2}{\hbar^2}
    }.
\end{equation*}

Так как
\begin{equation*}
    1+4\ell(\ell+1)
    =
    (2\ell+1)^2,
\end{equation*}
можно переписать
\begin{equation*}
    \tilde\ell
    =
    -\frac12
    +
    \frac12(2\ell+1)
    \sqrt{
        1
        -
        C_1
        \frac{8mZ_a e^2}
        {\hbar^2(2\ell+1)^2}
    }.
\end{equation*}

Считаем поправку малой и используем разложение
\begin{equation*}
    \sqrt{1-x}
    \simeq
    1-\frac{x}{2}.
\end{equation*}

Тогда
\begin{align}
    \tilde\ell
    &\simeq
    -\frac12
    +
    \frac12(2\ell+1)
    \left[
        1
        -
        C_1
        \frac{4mZ_a e^2}
        {\hbar^2(2\ell+1)^2}
    \right]
    \\
    &=
    \ell
    -
    C_1
    \frac{mZ_a e^2}
    {\hbar^2(\ell+1/2)}.
\end{align}

Введём величину
\begin{equation*}
    \boxed{
        \Delta_\ell
        =
        C_1
        \frac{mZ_a e^2}
        {\hbar^2(\ell+1/2)}
    }.
\end{equation*}

Тогда
\begin{equation*}
    \boxed{
        \tilde\ell
        =
        \ell-\Delta_\ell
    }.
\end{equation*}

Величина \(\Delta_\ell\) называется \textbf{квантовым дефектом}
или \textbf{поправкой Ридберга}.

\subsection{Эффективное главное квантовое число}

Для водородоподобной задачи главное квантовое число имеет вид
\begin{equation*}
    n=n_r+\ell+1.
\end{equation*}

После замены
\begin{equation*}
    \ell\rightarrow\tilde\ell
\end{equation*}
получаем
\begin{equation*}
    \tilde n
    =
    n_r+\tilde\ell+1.
\end{equation*}

Так как
\begin{equation*}
    \tilde\ell=\ell-\Delta_\ell,
\end{equation*}
то
\begin{align}
    \tilde n
    &=
    n_r+\ell+1-\Delta_\ell
    \\
    &=
    n-\Delta_\ell.
\end{align}

Таким образом,
\begin{equation*}
    \boxed{
        \tilde n=n-\Delta_\ell
    }.
\end{equation*}

Формула Бора для водородоподобного уровня
\begin{equation*}
    E_n
    =
    -\frac{Z_a^2Ry}{n^2}
\end{equation*}
переходит в
\begin{rbox}{[]}
\begin{equation*}
    \boxed{
        E_{n\ell}
        =
        -\frac{Z_a^2Ry}
        {(n-\Delta_\ell)^2}
    }.
\end{equation*}
\end{rbox}

Для нейтрального щелочного атома \(Z_a=1\):
\begin{equation*}
    \boxed{
        E_{n\ell}
        =
        -\frac{Ry}
        {(n-\Delta_\ell)^2}
    }.
\end{equation*}

\section{Физический смысл квантового дефекта}

Для атома водорода
\begin{equation*}
    E_n
    =
    -\frac{Ry}{n^2},
\end{equation*}
то есть энергия зависит только от главного квантового числа \(n\).

Поэтому состояния
\begin{equation*}
    ns,\qquad
    np,\qquad
    nd,\qquad
    nf,\ldots
\end{equation*}
с одинаковым значением \(n\) имеют одинаковую энергию.

Это называется случайным, или кулоновским, вырождением.

Для щелочного атома
\begin{equation*}
    E_{n\ell}
    =
    -\frac{Ry}{(n-\Delta_\ell)^2},
\end{equation*}
поэтому энергия уже зависит от \(\ell\).

Следовательно,
\begin{equation*}
    \boxed{
        \text{кулоновское вырождение по }\ell
        \text{ снимается}.
    }
\end{equation*}

Из выражения для квантового дефекта
\begin{equation*}
    \Delta_\ell
    \propto
    \frac{1}{\ell+1/2}
\end{equation*}
следует, что наиболее сильно возмущаются состояния с малыми
значениями \(\ell\).

Качественно
\begin{equation*}
    \Delta_s
    >
    \Delta_p
    >
    \Delta_d
    >
    \Delta_f
    >\ldots
\end{equation*}

При больших \(\ell\)
\begin{equation*}
    \Delta_\ell\rightarrow0.
\end{equation*}

Причина состоит в различной способности электронов проникать внутрь
атомного остова.

Для \(s\)-состояния
\begin{equation*}
    \ell=0,
\end{equation*}
поэтому центробежный барьер отсутствует, и электрон способен глубоко
проникать внутрь остова. Там экранирование ядра внутренними электронами
уменьшается, поэтому валентный электрон ощущает более сильное
кулоновское притяжение.

С ростом \(\ell\) увеличивается центробежный потенциал
\begin{equation*}
    U_{\text{цб}}(r)
    =
    \frac{\hbar^2\ell(\ell+1)}{2mr^2}.
\end{equation*}

Поэтому электрон всё хуже проникает внутрь атомного остова,
и его движение становится всё ближе к движению электрона в чистом
кулоновском поле.

При фиксированном \(n\) получаем порядок уровней
\begin{equation*}
    \boxed{
        E_{ns}
        <
        E_{np}
        <
        E_{nd}
        <
        E_{nf}
    }.
\end{equation*}

\section{Квантовые дефекты щелочных металлов}

Для основных состояний ряда щелочных металлов:

\begin{center}
\begin{tabular}{c|c|c}
    \hline
    Элемент & \(n\) & \(\Delta\) \\
    \hline
    Li & \(2\) & \(0.41\) \\
    Na & \(3\) & \(1.37\) \\
    K  & \(4\) & \(2.23\) \\
    Rb & \(5\) & \(3.20\) \\
    Cs & \(6\) & \(4.13\) \\
    \hline
\end{tabular}
\end{center}

Для данного элемента и фиксированного \(\ell\) квантовый дефект
слабо меняется при изменении \(n\), поэтому целую последовательность
возбуждённых уровней можно приближённо описывать одной формулой
\begin{equation*}
    E_{n\ell}
    =
    -\frac{Ry}{(n-\Delta_\ell)^2}.
\end{equation*}

% ============================================================
% РИСУНОК:
% График зависимости квантового дефекта Delta_l от l
% для Li, Na, K, Rb, Cs.
%
% По горизонтали l = 0,1,2,3.
% По вертикали Delta_l.
%
% Все кривые быстро стремятся к нулю с ростом l.
% Самые большие значения у Cs, затем Rb, K, Na, Li.
% ============================================================

\section{Энергетические уровни атома натрия}

Рассмотрим атом натрия
\begin{equation*}
    \mathrm{Na}:
    1s^2\,2s^2\,2p^6\,3s.
\end{equation*}

Основное состояние внешнего электрона соответствует уровню
\begin{equation*}
    3s.
\end{equation*}

Энергия ионизации натрия равна приблизительно
\begin{equation*}
    I_{\mathrm{Na}}
    \simeq
    5.13\ \text{эВ}.
\end{equation*}

Поэтому относительно границы ионизации
\begin{equation*}
    E_{3s}
    \simeq
    -5.13\ \text{эВ}.
\end{equation*}

Для следующих уровней из схемы получаем приблизительно
\begin{equation*}
    E_{3p}
    \simeq
    -2.0\ \text{эВ},
\end{equation*}
\begin{equation*}
    E_{3d}
    \simeq
    -1.6\ \text{эВ}.
\end{equation*}

Для сравнения, в атоме водорода все состояния с \(n=3\) имеют одну энергию:
\begin{equation*}
    E_3^{(H)}
    =
    -\frac{Ry}{3^2}
    \simeq
    -1.5\ \text{эВ}.
\end{equation*}

Таким образом, состояние \(3d\) натрия уже близко к водородоподобному
уровню, тогда как состояние \(3s\) сильно смещено вниз.

Это связано с тем, что \(s\)-электрон значительно глубже проникает
внутрь атомного остова.

% ============================================================
% РИСУНОК:
% Энергетические уровни Na и H.
%
% Слева четыре столбца S, P, D, F для натрия.
% Справа уровни водорода n=2,3,4,5,6,... и граница n=infinity.
%
% Для Na отметить:
% 3s ~ -5.13 эВ
% 3p ~ -2.0 эВ
% 3d ~ -1.6 эВ
%
% Показать, что при росте l уровни Na приближаются к
% соответствующим водородным уровням.
% ============================================================

\section{Диаграмма Гротриана}

Для систематизации спектральных уровней и переходов удобно использовать
диаграмму Гротриана.

Уровни располагают в столбцах в соответствии с орбитальным квантовым
числом:
\begin{equation*}
    \begin{array}{c|cccc}
        \ell & 0 & 1 & 2 & 3 \\
        \hline
        & S & P & D & F
    \end{array}.
\end{equation*}

Для электрических дипольных переходов действует правило отбора
\begin{equation*}
    \boxed{
        \Delta\ell=\pm1
    }.
\end{equation*}

Поэтому разрешены переходы
\begin{equation*}
    S\leftrightarrow P,
    \qquad
    P\leftrightarrow D,
    \qquad
    D\leftrightarrow F.
\end{equation*}

Например, прямой дипольный переход
\begin{equation*}
    S\leftrightarrow D
\end{equation*}
запрещён.

% ============================================================
% РИСУНОК:
% Диаграмма Гротриана для атома Na.
%
% Столбцы S, P, D, F.
% Основное состояние 3s.
% Показать уровни 3p, 4s, 4p, 3d, 4d, 4f, ...
% Показать разрешённые стрелками переходы Delta l = +/- 1.
%
% Подписать серии:
% главная,
% резкая,
% диффузная,
% фундаментальная.
% ============================================================

\subsection{Главная серия}

Главная серия соответствует переходам
\begin{equation*}
    nP\rightarrow3S.
\end{equation*}

Нижний уровень здесь является основным состоянием атома натрия.

Головная линия главной серии:
\begin{equation*}
    3P\rightarrow3S.
\end{equation*}

\subsection{Резкая серия}

Резкая серия соответствует переходам
\begin{equation*}
    nS\rightarrow3P.
\end{equation*}

\subsection{Диффузная серия}

Диффузная серия соответствует переходам
\begin{equation*}
    nD\rightarrow3P.
\end{equation*}

\subsection{Фундаментальная серия}

Фундаментальная серия соответствует переходам
\begin{equation*}
    nF\rightarrow3D.
\end{equation*}

Названия \emph{резкая} и \emph{диффузная} являются историческими
и связаны с внешним видом спектральных линий.

\section{Тонкая структура щелочного атома}

До сих пор спин-орбитальное взаимодействие не учитывалось.

Для внешнего электрона
\begin{equation*}
    s=\frac12.
\end{equation*}

Полный момент импульса
\begin{equation*}
    \mathbf J
    =
    \mathbf L+\mathbf S.
\end{equation*}

По правилу сложения моментов при \(\ell\neq0\)
\begin{equation*}
    \boxed{
        j
        =
        \ell\pm\frac12
    }.
\end{equation*}

Поэтому каждый уровень при \(\ell\neq0\) расщепляется на два подуровня.

Например,
\begin{equation*}
    P
    \longrightarrow
    P_{1/2},
    \ P_{3/2},
\end{equation*}
\begin{equation*}
    D
    \longrightarrow
    D_{3/2},
    \ D_{5/2},
\end{equation*}
\begin{equation*}
    F
    \longrightarrow
    F_{5/2},
    \ F_{7/2}.
\end{equation*}

Для \(S\)-состояния
\begin{equation*}
    \ell=0,
\end{equation*}
поэтому возможно только
\begin{equation*}
    \boxed{
        j=\frac12
    }.
\end{equation*}

Следовательно, \(S\)-терм на два значения \(j\) не расщепляется.

Таким образом:

\begin{itemize}
    \item вырождения по \(\ell\) в щелочном атоме уже нет;
    \item каждый уровень при \(\ell\neq0\) дополнительно расщепляется
    за счёт спин-орбитального взаимодействия.
\end{itemize}

% ============================================================
% РИСУНОК:
% Диаграмма Гротриана Na с тонкой структурой.
%
% Для P, D, F показать расщепление на два уровня:
% j = l-1/2 и j = l+1/2.
%
% Для S только j=1/2.
% ============================================================

\subsection{Поправка тонкой структуры}

Для водородоподобного уровня поправка тонкой структуры записывается
в виде
\begin{equation*}
    \boxed{
        \delta E_{\mathrm{fs}}
        =
        -\frac{\alpha^2Z_a^4Ry}{n^3}
        \left[
            \frac{1}{j+1/2}
            -
            \frac{3}{4n}
        \right]
    }.
\end{equation*}

Здесь
\begin{equation*}
    \alpha
    =
    \frac{e^2}{\hbar c}
    \simeq
    \frac{1}{137}
\end{equation*}
--- постоянная тонкой структуры.

Так как энергия зависит от \(j\), состояния
\begin{equation*}
    j=\ell-\frac12
\end{equation*}
и
\begin{equation*}
    j=\ell+\frac12
\end{equation*}
имеют немного различные энергии.

Поэтому спектральные линии могут расщепляться на дублеты.

Для главной линии натрия переход
\begin{equation*}
    3P\rightarrow3S
\end{equation*}
с учётом тонкой структуры превращается в два перехода:
\begin{equation*}
    3P_{1/2}
    \rightarrow
    3S_{1/2},
\end{equation*}
\begin{equation*}
    3P_{3/2}
    \rightarrow
    3S_{1/2}.
\end{equation*}

Поэтому знаменитая жёлтая линия натрия вблизи
\begin{equation*}
    \lambda\simeq589\ \text{нм}
\end{equation*}
при достаточном спектральном разрешении наблюдается как дублет.

\section{Рентгеновские спектры}

До сих пор рассматривались главным образом переходы внешнего электрона.

Теперь рассмотрим ситуацию, когда из атома удаляется внутренний электрон.
При этом на одной из внутренних оболочек возникает вакансия.

Электрон с более высокой оболочки может перейти на освободившийся уровень,
а разность энергий выделяется в виде рентгеновского фотона:
\begin{equation*}
    \hbar\omega
    =
    E_i-E_f.
\end{equation*}

Внутренние электронные оболочки обозначаются буквами
\begin{equation*}
    K,\quad L,\quad M,\quad N,\ldots
\end{equation*}
в соответствии с
\begin{equation*}
    K\leftrightarrow n=1,
\end{equation*}
\begin{equation*}
    L\leftrightarrow n=2,
\end{equation*}
\begin{equation*}
    M\leftrightarrow n=3,
\end{equation*}
\begin{equation*}
    N\leftrightarrow n=4.
\end{equation*}

Характерные энергии удаления различных электронов могут отличаться
на порядки. Для натрия валентный \(3s\)-электрон связан с энергией
порядка нескольких электрон-вольт, тогда как удаление \(1s\)-электрона
требует энергии порядка \(10^3\) эВ.

Поэтому переходы между внутренними оболочками лежат в рентгеновском
диапазоне.

% ============================================================
% РИСУНОК:
% Спектры рентгеновского излучения.
%
% Концентрические оболочки:
% K: n=1
% L: n=2
% M: n=3
% N: n=4
% далее n=5
%
% Показать переходы:
% L -> K : K_alpha
% M -> K : K_beta
% N -> K : K_gamma
%
% M -> L : L_alpha
% N -> L : L_beta
% и т.д.
% ============================================================

\subsection{Рентгеновские серии}

Если переход заканчивается на \(K\)-оболочке, соответствующие линии
образуют \(K\)-серию.

Переход
\begin{equation*}
    L\rightarrow K
\end{equation*}
обозначается
\begin{equation*}
    K_\alpha.
\end{equation*}

Переход
\begin{equation*}
    M\rightarrow K
\end{equation*}
обозначается
\begin{equation*}
    K_\beta.
\end{equation*}

Переход
\begin{equation*}
    N\rightarrow K
\end{equation*}
обозначается
\begin{equation*}
    K_\gamma.
\end{equation*}

Аналогично, переход
\begin{equation*}
    M\rightarrow L
\end{equation*}
даёт линию
\begin{equation*}
    L_\alpha.
\end{equation*}

\section{Экранирование внутренних электронов}

Для внутреннего электрона энергию приближённо можно записать
в водородоподобной форме, заменив заряд ядра \(Z\) эффективным зарядом
\begin{equation*}
    Z_{\mathrm{eff}}
    =
    Z-a_n,
\end{equation*}
где \(a_n\) --- постоянная экранирования для соответствующей оболочки.

Тогда
\begin{equation*}
    \boxed{
        E_n
        =
        -\frac{Ry}{n^2}
        (Z-a_n)^2
    }.
\end{equation*}

\section{Закон Мозли}

Пусть электрон переходит из состояния с главным квантовым числом
\(n_i\) в состояние с \(n_f\), причём
\begin{equation*}
    n_i>n_f.
\end{equation*}

Энергия испущенного фотона равна
\begin{equation*}
    \hbar\omega
    =
    E_i-E_f.
\end{equation*}

Для начального состояния
\begin{equation*}
    E_i
    =
    -Ry
    \frac{(Z-a_i)^2}{n_i^2},
\end{equation*}
а для конечного
\begin{equation*}
    E_f
    =
    -Ry
    \frac{(Z-a_f)^2}{n_f^2}.
\end{equation*}

Следовательно,
\begin{equation*}
    \boxed{
        \hbar\omega
        =
        Ry
        \left[
            \frac{(Z-a_f)^2}{n_f^2}
            -
            \frac{(Z-a_i)^2}{n_i^2}
        \right]
    }.
\end{equation*}

Для данной рентгеновской серии можно приближённо заменить
\begin{equation*}
    a_i\simeq a_f\simeq a,
\end{equation*}
где \(a\) --- средняя постоянная экранирования.

Тогда
\begin{equation*}
    \hbar\omega
    \simeq
    Ry(Z-a)^2
    \left(
        \frac{1}{n_f^2}
        -
        \frac{1}{n_i^2}
    \right).
\end{equation*}

Отсюда
\begin{rbox}{[]}
\begin{equation*}
    \boxed{
        \sqrt{\frac{\hbar\omega}{Ry}}
        =
        (Z-a)
        \sqrt{
            \frac{1}{n_f^2}
            -
            \frac{1}{n_i^2}
        }
    }.
\end{equation*}
\end{rbox}

Для фиксированной спектральной линии \(n_i\) и \(n_f\) постоянны,
поэтому
\begin{equation*}
    \boxed{
        \sqrt{\omega}
        \propto
        Z-a
    }.
\end{equation*}

Это соотношение называется \textbf{законом Мозли}.

Таким образом, квадратный корень из частоты характеристического
рентгеновского излучения линейно зависит от атомного номера элемента.

\subsection{Линия \(K_\alpha\)}

Для линии \(K_\alpha\)
\begin{equation*}
    n_i=2,
    \qquad
    n_f=1.
\end{equation*}

Поэтому
\begin{align}
    \hbar\omega_{K_\alpha}
    &=
    Ry(Z-a)^2
    \left(
        1-\frac14
    \right)
    \\
    &=
    \frac34
    Ry(Z-a)^2.
\end{align}

Следовательно,
\begin{equation*}
    \boxed{
        \hbar\omega_{K_\alpha}
        =
        \frac34
        Ry(Z-a)^2
    }.
\end{equation*}

Или
\begin{equation*}
    \boxed{
        \sqrt{
            \frac{\hbar\omega_{K_\alpha}}{Ry}
        }
        =
        \frac{\sqrt3}{2}(Z-a)
    }.
\end{equation*}

Для линии \(K_\alpha\) на лекции приводилось эмпирическое значение
\begin{equation*}
    a\simeq1.17,
    \qquad
    Z<30.
\end{equation*}

Для линии \(L_\alpha\)
\begin{equation*}
    a\simeq7.9,
    \qquad
    Z>62.
\end{equation*}

% ============================================================
% РИСУНОК:
% Закон Мозли.
%
% По горизонтали Z.
% По вертикали sqrt(E/Ry) или sqrt(omega).
%
% Прямая линия с ненулевым пересечением оси Z,
% соответствующим постоянной экранирования a.
% ============================================================

\section{Итоговая структура уровней}

Для чистого кулоновского поля энергия зависит только от главного
квантового числа:
\begin{equation*}
    E=E(n).
\end{equation*}

Учёт взаимодействия валентного электрона с атомным остовом приводит
к появлению квантового дефекта:
\begin{equation*}
    E=E(n,\ell).
\end{equation*}

При этом снимается случайное кулоновское вырождение по \(\ell\):
\begin{equation*}
    E_{n\ell}
    =
    -\frac{Z_a^2Ry}
    {(n-\Delta_\ell)^2}.
\end{equation*}

После учёта спин-орбитального взаимодействия энергия зависит также
от полного момента электрона:
\begin{equation*}
    E=E(n,\ell,j).
\end{equation*}

Таким образом, последовательность уточнений имеет вид
\begin{equation*}
    \boxed{
        E(n)
        \longrightarrow
        E(n,\ell)
        \longrightarrow
        E(n,\ell,j)
    }.
\end{equation*}

Для внутренних электронов удобно пользоваться моделью экранированного
заряда
\begin{equation*}
    E_n
    =
    -\frac{Ry}{n^2}(Z-a_n)^2,
\end{equation*}
которая приводит к закону Мозли для характеристического
рентгеновского спектра.

\end{document}
